\section{总结}
\label{sec:summary}

本工作定义了一种从期望概率分布中抽取格点场构型样本的基于流的 MCMC 算法：
\begin{enumerate}
    \item 训练一个 real NVP 流模型，使其近似产生期望分布；
    \item 由训练好的模型提出样本；
    \item 从任意初始构型出发，用 Metropolis-Hastings 算法对每个提议做接受或拒绝以推进马尔可夫链。
\end{enumerate}
本文证明了该方案定义一条遍历且平衡的马尔可夫链，从而保证在长链极限下收敛到期望概率分布。本质上，基于流的 MCMC 算法把基于神经网络的归一化流的表达能力与马尔可夫链的理论保证结合起来，构造出一个可训练且渐近正确的采样器。由于这些流适用于任何具有连续实值自由度的构型，可以把该方法一般性地应用于一大类格点理论中的任意一个。这里针对 $\phi^4$ 理论具体实现了该算法，并通过若干物理观测量的研究表明，其产生的构型系综与用标准局域 Metropolis 和 HMC 算法生成的系综不可区分。

该方案的一个关键特点是：训练到固定接受率的模型在抽样阶段不会经历临界慢化。特别地，所有观测量的自相关时间完全由流模型训练的精度决定；完美训练对应去关联样本以及 MCMC 过程中 Metropolis-Hastings 步骤 100\% 的接受率。尽管如此，这一方案的训练环节能否高效扩展到更大的模型规模、以及 QCD 等更复杂的理论，仍有待研究。流模型训练与规模扩展方面的最新进展为这方面提供了乐观的理由。此外，引入对称性通常能改善训练的数据效率，实现时空对称性与规范对称性~\cite{Cohen2019GaugeEquiv} 可能是使这些流模型能在 QCD 这类格点规范理论上切实训练起来的自然下一步。

在走向 QCD 这类格点规范理论的过程中，还需要若干理论上的发展。此处选来参数化归一化流的 real NVP 模型是用变量向量 $\phi \in \mathbb{R}^D$ 描述的。然而规范组态生活在由李群结构产生的紧致流形上。扩展该方法将需要一个能作用于该流形、同时保持充分表达能力的归一化流模型。先验的选择同样需要推广为格点规范组态流形上一个易于抽样的分布。例如均匀分布可能是先验的一个候选，但这一选择必须在具体流模型的语境下检验。

如果本文提出的基于流的 MCMC 算法能够为 QCD 这样的复杂理论实现，其优势将是显著的：可以以极小的代价生成任意大的场构型系综。提议步骤独立于任何先前构型，允许并行生成提议；神经网络执行的硬件与软件支持持续改善，预示着这类系综生成方式未来的实用收益。既然能从训练好的模型高效生成样本，系综便无需长期存储。此外，为某个作用量训练的模型既可以重新训练，也可以作为针对参数值相近的另一作用量的流模型的先验。这将使参数调节更加高效，并能在现有模型之间插值或从其外推，生成更多系综。
